\chapter{Introduction}
\label{chap:introduction}


% use [] to set name for ToC
\section[Context \& Motivation]{Context \& Motivation} % ok with fontsize=12pt
    Open-source software is the foundation of modern computing. 
    Individual users and companies alike rely on openly available source code, maintained by
    communities of volunteers and by companies alike.
    Open-source has become effectively universal: about 98\% of the audited codebases contain open source components~\cite{ossra2026}.

    Open-source security is a double-edged sword: while extreme transparency enables faster bug detection and patching by the community,
     it also provides attackers with everything they needed to hunt for vulnerabilities.
    In addition, Open-source dependencies are nested in the majority of software projects, leading to possible supply-chain attacks.

    The most important and known service in term of OSS\footnote{OSS stands for open-source software.} is GitHub: a global platform for software development based on version control and collaboration.
    It hosts the vast majority of the open source software in use today, with more than two hundred million active repositories~\cite{ayala2025disclosures}. 
    At the base of the platform there is the concept of commit: a snapshot of the source code at a specific point of time.
    The direct consequence of the structure of this framework is the possibility of inspecting any change made to the repository as soon as it is published.
    The nature of the commit can be various, from bug correction to code refactoring. For the scope of this thesis, the attention will be posed on secuirity updates.

    Security fixes, however, do not reach the public in the same way as ordinary changes.
    Under the CVD\footnote{CVD stands for Coordinated Vulnerability Disclosure.} model, a vulnerability is corrected before it is
    announced, so that a remedy already exists by the time users learn that they are at risk.
    During this embargo period, between the correction and the disclousure of a vulnerability, the fix is committed like any other change, but the commit message is
    kept vague or generic, so that the change does not attract the attention of anyone watching the
    repository~\cite{colefunda2023}.
    

    Between the moment a fix is written and the moment it actually protects there are two
    distinct delays~\cite{imtiaz2023opensneaky}.
    The first one lies between the commit and the release: the correction is already in the
    repository, but it only becomes installable when the maintainers publish a new version of the
    package.
    The second one lies between that release and the publication of a security advisory, which is
    what tells the automated tools used by client projects that they should update.
    In both cases the code of the fix is already public, while the projects that depend on it are
    still exposed: during the first delay users cannot update at all, because no release exists yet,
    and during the second one they could, but nothing tells them that they should.

    This practice is known as silent patch: the vulnerability is fixed, but nothing in the commit
    says that a vulnerability was ever there.
    A direct consequence is that users keep running vulnerable software, because they receive no
    signal telling them that a specific update matters for the security.

    Outdated code is the norm rather than the exception: one study of more than 4,600 GitHub
    projects found that 81.5\% of them kept dependencies that were no longer up to
    date~\cite{ayala2025disclosures}, and about 92\% of the audited codebases contain components
    that are ten or more versions behind the current release~\cite{ossra2026}.

    Keeping the commit message quiet, however, does not hide the vulnerability, because the
    information that describes it is the change itself.
    A silent patch shows which lines were wrong and how they were repaired, so anyone who reads it
    can work backwards from the correction to the flaw corrected.
    In addition, the openness of the code allows to search for weaknesses that nobody has reported and nobody has fixed 
    even without the presence of a security patch.

    Artificial intelligence is a relatively new technology that is developing quickly, and it has
    deeply changed cybersecurity, creating new scenarios, new kinds of attack, new ways of
    defending against them, and new ways to analyse code in search of vulnerabilities.

    
    The analysis of unfamiliar code to find security problems is an expensive activity
    in terms of time and knowledge required. As a consequence, the possibility for any user to use AI general purpose
    models could be a game-changer.

    Measuring how well an automated analysis performs on this task is harder than it seems, and the
    first obstacle is the data.
    What is needed is a set of vulnerabilities paired with the commit that fixes each of them, but
    this pairing is rarely available: an empirical study of 193,448 \footnote{CVE stands for Common Vulnerabilities and Exposures} entries found that only
    30.99\% of them include an explicit link to a patch commit, and that 41.34\% of those links no
    longer work, because repositories get deleted, branches are renamed or pull requests are
    merged~\cite{gitpatchdb2026}.
    Furthermore, in many cases, commits do not provide any help to the data gathering phase since we are
     talking about silent patches.
    In addition, vulnerabilities often emerge through chains of calls that cross several
    functions rather than within isolated ones~\cite{yildiz2025jitvul}, 
    and the fixes that touch more than one file are exactly the ones where approaches 
    working at function level struggle.
    Finally, saying that a vulnerability exists is not enough, because it also has to be classified
    with the corresponding CWE\footnote{CWE stands for Common Weakness Enumeration}, chosen among hundreds of categories 
    that could differ from one another by subtle distinctions.

    The problem is not new, and neither are the attempts to automate it.
    A first generation of work applied deep learning to source code in order to detect
    vulnerabilities~\cite{vuldeepecker2018}, but a later evaluation, under realistic conditions,
    found that the performance of such models dropped~\cite{reveal2022}.
    A parallel line of research addressed the narrower task of recognising silent fixes among
    ordinary commits~\cite{colefunda2023}.

    The arrival of Large Language Models reopened the question.
    When each vulnerable function is paired with its corrected version, models (used on their
    own) label almost everything as vulnerable, while agents that explore the repository by
    themselves are better at discriminating between the two versions~\cite{yildiz2025jitvul}.
    Agents that are also allowed to compile, execute and fuzz the code have already found real
    vulnerabilities in open source projects~\cite{fuzzingbrain2026}, and yet, when they are asked
    to find and fix a vulnerability that is genuinely new, even agents are not able to solve it on their own~\cite{zerodaybench2026}.
    

\section[Methodology]{Methodology}

    This work aims to answer three questions that follow from the situation described above.
    Firstly, whether an agent running on a general purpose model is able to identify and classify the
    vulnerability when it is given only the code before and after the commit and the diff file
    containing the changes.
    Furthermore, how this ability to detect and classify the vulnerability changes when more information is
    added to the context.
    Finally, whether the agent is able to detect and classify a vulnerability with access to the
    repository alone and no further information, which simulates the search for a zero-day vulnerability.

    The proposal is an architecture based on AI agents, evaluated in three versions that differ
    only in what the agent is allowed to see:
    \begin{itemize}
        \item the agent receives only the code changed by a commit, in its version before and
              after the correction together with the corresponding diff;
        \item it receives the whole repository together with those same changes;
        \item it receives only the repository, with no indication that a fix exists or where to
              look.
    \end{itemize}

    Each version is asked for the same answer: whether a vulnerability is present and which
    CWE it corresponds to, so that the three results can be compared directly.

    As a secondary contribution, the thesis also lays the groundwork for a real time implementation.
    The evaluation above needs a vulnerability that is already known, which is what makes the
    measurement possible but also confines it to the past.
    The tool developed for this purpose works in the opposite direction: it selects the
    repositories that are worth keeping an eye on for a future real-time analysis by the described architecture,
     and then follows them over time, recording which
    of the selected projects end up with a published CVE and how many days pass between the moment
    they were selected and that publication.
    The results it produces are preliminary and are kept separate from the evaluation of the
    agents, which does not rely on them in any way.
    


\section[Thesis Organization]{Thesis Organization}

    This thesis is structured in seven chapters:

    \begin{itemize}
        \item \textbf{Chapter 2 -- Background} introduces the concepts the work relies on. It
              describes how vulnerabilities are disclosed and catalogued, what makes a patch
              silent, and how large language models are turned into agents that reason and use
              tools.

        \item \textbf{Chapter 3 -- Related Work} reviews the main lines of research on the
              automated detection of vulnerabilities, including the approaches that rely on
              language models and agents, and positions the present work with respect to them.

        \item \textbf{Chapter 4 -- Methodology} details the data collection and the proposed
              architecture. It explains how the dataset used for the evaluation was gathered and
              built, and then presents the three agent configurations, the isolation under which
              each run is executed and the format in which a verdict has to be returned.

        \item \textbf{Chapter 5 -- Repository Monitoring} presents the tool developed as a
              secondary contribution. It explains how the repositories worth watching are
              selected, and how the selection is verified over time against the vulnerabilities
              that are published afterwards.

        \item \textbf{Chapter 6 -- Results} reports the experimental results. It defines the
              evaluation metrics, compares the three configurations and the models under test, and
              analyses the cases in which the agents failed.

        \item \textbf{Chapter 7 -- Conclusions} summarises the findings and discusses directions
              for future work, including the real time extension built on the monitoring tool.
    \end{itemize}
